\documentclass[noproofs]{rmj-public}
\usepackage{amsmath,amssymb,booktabs,array,needspace}
\usepackage{xurl}
\newtheorem{thm}{Theorem}[section]
\newtheorem{theorem}[thm]{Theorem}
\newtheorem{lemma}[thm]{Lemma}
\newtheorem{proposition}[thm]{Proposition}
\newtheorem{corollary}[thm]{Corollary}
\theoremstyle{definition}
\newtheorem{definition}[thm]{Definition}
\newtheorem{remark}[thm]{Remark}
\numberwithin{equation}{section}
\newcommand{\R}{\mathbb R}
\newcommand{\Q}{\mathbb Q}
\newcommand{\Gal}{\operatorname{Gal}}
\newcommand{\RMC}{\mathrm{RMC}}
\newcommand{\code}[1]{\texttt{\small\detokenize{#1}}}
% Suppress production-only volume, issue, and DOI placeholders in a submission.
% The official class and its page dimensions are unchanged.
\renewcommand{\logobox}{}
\title[Real cyclic quintic extensions]{Restricted marked-ruler-and-compass constructibility of real cyclic quintic extensions}
\author{}
% No author is assigned in this copy.
\date{October 16, 2026}
\subjclass[2020]{Primary 51M15; Secondary 12F10}
\keywords{Geometric construction, marked ruler, cyclic quintic extension, verging polynomial, formal verification}
\hypersetup{pdftitle={Restricted marked-ruler-and-compass constructibility of real cyclic quintic extensions},pdfauthor={},hidelinks}

\begin{document}
\begin{abstract}
We prove that every element of a real cyclic extension of degree five is
constructible relative to its base field by restricted marked ruler and
compass. A universal identity in adjacent products of five conjugates
produces an auxiliary cubic and a line--circle verging sextic. Strict
admissibility inequalities on a rational template persist on a common
neighborhood of a compact interval. Density of realizable conjugate tuples
then supplies suitable data in every real cyclic quintic extension.
Two adjacent-product values recover a generator of the original extension
by field operations. The coordinate-constructibility theorem is verified
in Lean, and a written soundness argument translates its finite rules into
the permitted geometric operations. An exact conductor-25 example and a
sufficient condition for regular-polygon constructibility are also given.
\end{abstract}
\maketitle

\section{Introduction and construction model}\label{sec:intro}

A marked ruler permits a line to be positioned through a fixed point with
two specified marks on prescribed curves. We consider a ruler with marks
one unit apart and the restricted marked-ruler-and-compass model of
Benjamin and Snyder \cite{BS14,BS16}. In this model verging is permitted
between two lines or between a line and a circle. Their construction of the
regular hendecagon \cite{BS14} provides a cyclic quintic example. Their later
work \cite{BS16} treats certain fifth roots of cubic irrationals. The result
here concerns an arbitrary real cyclic extension of degree five, with a
uniform existence argument based on a rational admissible template.

\begin{definition}\label{def:rmc}
Let $K\subseteq\R$ be a subfield. A real number is \emph{RMC-constructible
relative to $K$} if its point on the horizontal axis occurs in a finite
construction sequence starting from $(0,0)$, $(1,0)$ and available points
$(k,0)$, $k\in K$. A step is an ordinary straightedge-and-compass
construction, or the selection of a permissible pair of unit-separated
ruler marks on two previously specified lines, or on a previously
specified line and circle, with their ruler line passing through an
already constructed point.
\end{definition}

Every individual sequence uses only finitely many initial field elements.
When $K$ is generated over $\Q$ by finitely many supplied coefficients,
each required element of $K$ is obtained from them by field operations.
The definition is existential: a specified permissible placement may be
selected. It does not assert that an arbitrary placement produces the
desired root. There is no betweenness requirement on the fixed point.

\begin{theorem}\label{thm:main}
Let $K\subseteq\R$ and let $L/K$ be a Galois extension of degree five
embedded in $\R$. Every element of $L$ is RMC-constructible relative to $K$.
\end{theorem}

The Galois group has prime order five and is cyclic. Conversely, an
irreducible quintic over a real field with Galois group $C_5$ has only real
roots: complex conjugation on its splitting field has order dividing two,
and hence is trivial. Thus the theorem covers all its roots, relative to
the field generated by the input coefficients.

The proof first constructs a real auxiliary $u$ of degree at most three
over $K$. Three positive square roots over $K(u)$ prepare a line and a
circle. Two selected verging outputs then recover a generator of $L$.
The square roots need not belong to $K(u)$; their constructions are
included. No assertion of a minimal operation count is made.

The central coordinate theorem and its algebraic, topological, and
field-theoretic ingredients are checked in the accompanying Lean source.
The interpretation of its rules as geometric constructions is proved
below. The explicit example and cyclotomic corollary are written and
computational supplements; they are not separate theorems in that source.

\section{Universal adjacent-product identities}\label{sec:algebra}

Indices in this section are taken modulo five. Initially
$x_0,\ldots,x_4,U$ are independent variables over $\Q$. Define
\begin{equation}\label{eq:g}
g_i=1+Ux_i,\qquad \beta_i=g_ig_{i+1},\qquad
P=\prod_{i=0}^4g_i,\qquad T=\sum_{i=0}^4\beta_i,\qquad D=T+1.
\end{equation}
Write $e_j=e_j(\beta_0,\ldots,\beta_4)$, $e_0=1$, and
$R=\sum_i g_ig_{i+1}g_{i+2}$. The cycle gives
\begin{equation}\label{eq:products}
e_5=P^2,\qquad e_4=PR.
\end{equation}
Indeed, the product of all five adjacent products is $P^2$. The product
omitting a single adjacent factor is $P$ times the product of the three
complementary $g$ factors, giving the second identity on summation.

For polynomials in the original $x$ variables, put
\begin{align*}
S&=\sum_i x_i,& t&=\sum_i x_ix_{i+1},& E_j&=e_j(x_0,\ldots,x_4),\\
v&=\sum_i x_ix_{i+1}^2x_{i+2},& w&=\sum_i x_ix_{i+1}x_{i+2},\\
A_{32}&=\sum_i(x_i^3x_{i+1}^2+x_i^2x_{i+1}^3),&
A_{33}&=\sum_i x_i^3x_{i+1}^3.
\end{align*}
Define
\begin{align}\label{eq:C}
C_0&=4(t^2-4v),\notag\\
C_1&=4A_{32}-2St^2+2t(w+E_3)+4SE_4+4E_5,\notag\\
C_2&=(4A_{33}-t^3)/3+2tE_4+4SE_5,\notag\\
C_3&=2tE_5,\qquad H_x(U)=C_0+C_1U+C_2U^2+C_3U^3.
\end{align}
Each $C_j$ is invariant under cyclic rotation of the $x$ variables.

\begin{lemma}\label{lem:universal}
The following polynomial identity holds:
\begin{equation}\label{eq:universal}
\begin{split}
D^3-4(e_2+T)D+4(e_3+e_2)
+2(P+R)D-8P=U^4H_x(U).
\end{split}
\end{equation}
\end{lemma}
\begin{proof}
Newton's identities rewrite the left side as
\[
\frac{4\sum_i\beta_i^3-T^3}{3}-T^2-T+1
+2(T+1)R+2(T-3)P.
\]
Substitute
\begin{align*}
T&=5+2SU+tU^2,\\
R&=5+3SU+(t+E_2)U^2+wU^3,\\
P&=1+E_1U+E_2U^2+E_3U^3+E_4U^4+E_5U^5,\\
\sum_i\beta_i^3&=\sum_i(1+Ux_i)^3(1+Ux_{i+1})^3.
\end{align*}
The coefficients through degree three cancel. The coefficients in degrees
four through seven are respectively the four expressions in
\eqref{eq:C}. This proves the identity by coefficient comparison.
\end{proof}

Assume $D\ne0$ and set
\begin{equation}\label{eq:sextic}
\lambda=\frac2D,\qquad
f(z)=(z-\lambda)\prod_{i=0}^4(z-\lambda\beta_i)
=z^6-2z^5+\sum_{j=2}^6a_jz^{6-j}.
\end{equation}
Since $\lambda(1+T)=2$, its coefficients satisfy
\begin{equation}\label{eq:aj}
a_j=(-\lambda)^j(e_j+e_{j-1})\quad(2\le j\le5),
\qquad a_6=\lambda^6P^2.
\end{equation}
Keep track of the sign by defining
\begin{equation}\label{eq:h}
h=-\lambda^3P,\qquad m=2-2a_2-a_3.
\end{equation}
In particular $a_6=h^2$.

\begin{lemma}\label{lem:signed}
Under the preceding definitions,
\begin{equation}\label{eq:signed}
a_5+mh+2h^2=-\frac{16PU^4}{D^6}H_x(U).
\end{equation}
\end{lemma}
\begin{proof}
Substitution of \eqref{eq:aj} and \eqref{eq:h} gives
\begin{align*}
D^6(a_5+mh+2h^2)
={}&-32(P^2+e_4)D-16PD^3\\
&+64P(e_2+T)D-64P(e_3+e_2)+128P^2.
\end{align*}
Use $e_4=PR$ to identify this expression with $-16P$ times
the left side of \eqref{eq:universal}, and apply that identity.
\end{proof}

\section{The permitted geometry}\label{sec:geometry}

\subsection{Soundness of the coordinate rules}
The coordinate predicate used in Lean starts with $K$ and is inductively
closed under addition, negation, multiplication, inverse, nonnegative
square root, and the following two rules. Given available $a,b,c,r$ with
$a\ne0$, $r>0$, either $X$ or $Y$ may be constructed if there are real
$X,Y,V,W$ such that
\begin{equation}\label{eq:constraints}
\begin{gathered}
V=a,\qquad (X-b)^2+(Y-c)^2=r^2,\\
XW=YV,\qquad (X-V)^2+(Y-W)^2=1.
\end{gathered}
\end{equation}

\begin{lemma}[Geometric soundness]\label{lem:sound}
Every finite derivation in these coordinate rules is realizable by a
finite RMC construction relative to $K$.
\end{lemma}
\begin{proof}
Induct on the derivation. Field arithmetic and nonnegative square roots
are ordinary straightedge-and-compass constructions with directed axis
segments. For example, the line through $(0,0),(1,y)$ meets the vertical
line of first coordinate $x$ at $(x,xy)$. For $x\ne0$, the line through
$(0,0),(x,1)$ meets the vertical line of first coordinate 1 at $(1,1/x)$.
In Lean $0^{-1}=0$; this case reuses the origin. If $t>0$, the circle
with diameter $[-1,t]$ on the horizontal axis meets the positive vertical
axis at height $\sqrt t$. If $t=0$, reuse the origin. Directed addition,
reflection, coordinate pairs, perpendicular projections and quarter-turns
are ordinary finite compass constructions.

For \eqref{eq:constraints}, construct the line through $(a,0),(a,1)$
and the circle centered at $(b,c)$ through $(b+r,c)$. Each uses distinct
constructed points, and the line avoids the fixed origin. Write
$P=(X,Y)$, $Q=(V,W)$. Then $Q$ is on the line, $P$ on the circle,
and their distance is one. Since $Q\ne0$, the determinant equation says
that $O,P,Q$ are collinear. If $P\ne O$, these are precisely the
constraints for an allowed placement. Its primary point supplies both
coordinates by projections. If $P=O$, both outputs are already zero.
Only $a,b,c,r$ are required as earlier inputs; the other coordinates are
outputs of the placement. Finally, concatenating finitely many premise
constructions and the new finite construction preserves finiteness.
\end{proof}

\subsection{Direct reconstruction from the sextic}

\begin{lemma}\label{lem:reconstruct}
Suppose
$f(z)=z^6-2z^5+a_2z^4+a_3z^3+a_4z^2+a_5z+h^2$,
where $m=2-2a_2-a_3$ and $a_5=-mh-2h^2$. Define
\begin{equation}\label{eq:parameters}
q=\frac{a_3-m-4h}{8},\qquad
A=\frac{a_4-(h+m/2)^2+2h+4q}{4q},\qquad
R_0=h+\frac{m^2}{16A}+q.
\end{equation}
If $q,A,R_0>0$, put
\[
a=\sqrt A,\quad c=\sqrt q,\quad b=\frac m{4a},\quad
\rho=\sqrt{R_0}.
\]
Every real root $z\ne0,1$ of $f$ is recovered by field operations from a
permissible unit-marked-ruler placement through the origin, between
the line $X=a$ and the circle centered at $(b,c)$ of radius $\rho$.
If $h\ne0$, the excluded values are not roots.
\end{lemma}
\begin{proof}
We have $4ab=m$ and $\rho^2=b^2+c^2+h$. Expansion gives
\begin{equation}\label{eq:geompoly}
f(z)=\big((z^2-h)(z-1)-2abz\big)^2
+4a^2c^2z^2-4c^2z^2(z-1)^2.
\end{equation}
For instance, writing $k=h+m/2$, the five coefficients after $-2$ are
\[
1-2k-4q,\quad 2h+2k+8q,\quad
k^2-2h+4Aq-4q,\quad -2hk,\quad h^2,
\]
which equal $a_2,a_3,a_4,a_5,h^2$ by \eqref{eq:parameters}.
For a root $z\ne0,1$, set
\begin{equation}\label{eq:marks}
X=\frac{az}{z-1},\qquad
Y=\frac{z^2-h-2bX}{2c},\qquad P_0=(X,Y),\qquad
Q_0=\frac{z-1}{z}P_0.
\end{equation}
Equation \eqref{eq:geompoly} implies $X^2+Y^2=z^2$, and hence
\[
(X-b)^2+(Y-c)^2=\rho^2,\qquad
(Q_0)_x=a,\qquad \det(P_0,Q_0)=0,\qquad
|P_0-Q_0|^2=1.
\]
Thus \eqref{eq:marks} exhibits an actual placement, not merely a root of
a necessary equation obtained by elimination. The line and circle depend
only on the available coefficients; $z$ is used in the existence witness,
not as a geometric input. After selecting this placement,
\[
X-a=\frac a{z-1}\ne0,\qquad z=\frac X{X-a}.
\]
Finally $f(0)=h^2>0$ if $h\ne0$, while
$f(1)=4a^2(b^2+c^2)>0$.
\end{proof}

The coefficient formulas are consistent with the verging-polynomial
formulas of \cite[Theorem 3.1]{BS16}. The preceding witness argument
establishes the sufficiency needed here directly, including the sign and
denominator exceptions.

\subsection{Cubic construction without an algebraic-root rule}

\begin{lemma}\label{lem:cubic}
Every real root of a cubic with available coefficients is obtainable by
the coordinate rules above and hence by RMC construction.
\end{lemma}
\begin{proof}
Divide by the nonzero leading coefficient and translate to a depressed
cubic $u^3+pu+q=0$. If $q=0$, either $u=0$ or $u^2=-p$, and a square
root with the selected sign suffices.

Suppose $p,q\ne0$. Choose a positive integer $N$ large enough that
$v=u/N$, $p'=p/N^2$, $q'=q/N^3$ satisfy $|v|<1$ and $p'>-1/2$.
This can be chosen from coefficient data: $|u|\le1+|p|+|q|$, as follows
by taking absolute values in the equation when $|u|>1$. Put
\[
a=\frac{2q'}{p'},\qquad b=-\frac a{p'},\qquad
h=1+\frac{2a^2}{p'},\qquad r=\sqrt{b^2+h},\qquad X=v+a.
\]
The inputs are available, $a\ne0$, and
\[
b^2+h=1+\frac{a^2(1+2p')}{p'^2}>0.
\]
If $X=0$, recover $u=-Na$. Otherwise $v\ne0$ and the real witness
$Y=\sqrt{X^2(1/v^2-1)}$ exists. The identities
\begin{align*}
X^2+Y^2&=X^2/v^2,\\
v^2(h+2bX)-X^2&=(-2a/p')(v^3+p'v+q')=0
\end{align*}
show that $P=(X,Y)$ is on the circle centered at $(b,0)$ of radius $r$.
For $Q=(a/X)P$, its first coordinate is $a$, it is collinear with $P,O$,
and $|P-Q|^2=(1-a/X)^2(X^2+Y^2)=1$ because $X-a=v$. Thus the
permitted placement supplies $X$, and $u=N(X-a)$. The witness $Y$
is not assumed available before verging. A circle centered on the
horizontal axis is permitted by the construction model.

If $p=0$, $q\ne0$, set $y=u^2+u$. Direct expansion gives
\[
y^3+3qy+(q-q^2)=0,\qquad
u=\frac{y-q}{y+1},\qquad y+1=(u+1/2)^2+3/4>0.
\]
The auxiliary equation has nonzero linear coefficient and is covered by
the preceding case. Field operations then recover $u$.
\end{proof}

\section{An admissible rational template}\label{sec:template}

Take $x^*=(-3,-1,1,4,2)$. Equations \eqref{eq:C} give
\[
(C_0,C_1,C_2,C_3)=(912,792,792,384),\qquad
H_{x^*}(U)=24F(U),
\]
where $F(U)=16U^3+33U^2+33U+38$. Define
\[
\ell=-\frac{167710}{100000},\qquad
\eta=-\frac{167709}{100000}.
\]
Exact evaluation yields
\[
F(\ell)=-\frac{8943193}{31250000000}<0,
\qquad F(\eta)=\frac{17937874421}{62500000000000}>0.
\]
Also $F'(U)=48(U+11/16)^2+165/16>0$. There is therefore a unique real
template root in $(\ell,\eta)$, although uniqueness will not be needed.

\begin{lemma}\label{lem:template}
For every $U\in[\ell,\eta]$, the successive rational expressions above
are defined and satisfy
\[
D(x^*,U)>0,\quad P(x^*,U)<0,\quad
q(x^*,U)>0,\quad A(x^*,U)>0,\quad R_0(x^*,U)>0.
\]
\end{lemma}
\begin{proof}
Rational interval evaluation gives the enclosures in
Table~\ref{tab:template}. Arithmetic is performed successively, so each
denominator interval is first checked to exclude zero. The table implies
the stated signs. These strict signs throughout the interval are also
proved by rational interval certificates in the accompanying Lean source.
\end{proof}

\begin{table}[ht]
\centering
\caption{Exact enclosures throughout the template interval.}\label{tab:template}
\begin{tabular}{lrr}
\toprule
Quantity&Lower bound&Upper bound\\\midrule
$D$&$92191/5000$&$18439/1000$\\
$P$&$-734609/5000$&$-36729/250$\\
$\lambda$&$271/2500$&$217/2000$\\
$h$&$937/5000$&$469/2500$\\
$m$&$4801/5000$&$9621/10000$\\
$q$&$779/2000$&$39/100$\\
$A$&$1979/10000$&$2001/10000$\\
$R_0$&$173/200$&$8697/10000$\\\bottomrule
\end{tabular}
\end{table}

The interval bounds in Table~\ref{tab:template} are supplied by the
standard-library rational-arithmetic certificate; the Lean proof may use
different intermediate enclosures with the same strict-sign conclusions.
For specialization we also note $C_3(x^*)=384>0$ and $x^*_0<x^*_1$.
Pairwise distinctness of the adjacent products is not required.

\section{Simultaneous specialization}\label{sec:specialize}

\begin{lemma}\label{lem:specialize}
For every real cyclic quintic extension $L/K$ with generator $\sigma$,
there exist $\gamma\in L$ and $u\in(\ell,\eta)$ such that, for
$x_i=\sigma^i(\gamma)$, one has $H_x(u)=0$, $C_3(x)>0$, $x_0<x_1$,
and all five strict inequalities of Lemma~\ref{lem:template} at $(x,u)$.
\end{lemma}
\begin{proof}
Let $G\subseteq\R^5\times\R$ be the set of pairs with
$D>0$, $P<0$, $q>0$, $A>0$, $R_0>0$. It is open: polynomial
expressions are continuous, rational expressions are continuous where
their denominators are nonzero, and the defining signs guarantee the
successive denominator restrictions. By Lemma~\ref{lem:template},
$\{x^*\}\times[\ell,\eta]\subseteq G$.

For each $U$ in the interval choose a product neighborhood
$V_U\times J_U\subseteq G$ of $(x^*,U)$. Compactness supplies a finite
subcover by the $J_U$. Intersect the corresponding $V_U$ to obtain an
open neighborhood $V$ of $x^*$ with
\begin{equation}\label{eq:tube}
V\times[\ell,\eta]\subseteq G.
\end{equation}
Shrink $V$ using continuity of the endpoint evaluations, $C_3$ and the
coordinate difference, so that every $x\in V$ also satisfies
\[
H_x(\ell)<0<H_x(\eta),\qquad C_3(x)>0,\qquad x_0<x_1.
\]
The intermediate value theorem supplies a root $u\in(\ell,\eta)$
for each such $x$; \eqref{eq:tube} guarantees all admissibility
conditions at this same pair. No continuous choice of roots is needed.

Choose a $K$-basis $b_0,\ldots,b_4$ of $L$. The matrix
$B_{ij}=\sigma^i(b_j)$ is nonsingular. Indeed a linear relation over
$L$ among its rows would vanish on a $K$-basis and hence on every
element of $L$, contradicting independence of distinct field
homomorphisms. Since $L\subseteq\R$, this is a real invertible matrix.
Thus $B(\Q^5)$ is dense in $\R^5$. Choose $r\in\Q^5$ with $Br\in V$
and set $\gamma=\sum_jr_jb_j$. Its conjugate tuple is $Br$ because
$\sigma$ fixes $\Q$. Apply the preceding paragraph to this tuple.
\end{proof}

\section{Descent and recovery of the quintic generator}\label{sec:recovery}

\begin{proof}[Proof of Theorem~\ref{thm:main}]
Choose $\gamma,u$ by Lemma~\ref{lem:specialize}, and put
$K'=K(u)$. Rotation of the conjugate tuple is the action of $\sigma$.
Consequently each $C_j(x)$ is fixed by the Galois group and lies in $K$.
Since $C_3(x)>0$, $H_x$ is a nonzero cubic over $K$, so $[K':K]\le3$.
Lemma~\ref{lem:cubic} constructs its selected real root $u$.

The same invariance applied coefficientwise in $U$ shows that
$P,T,R,e_j$ belong to $K[U]$. Hence their values at $u$, and
$\lambda,a_j,h,m,q,A,R_0$, are available from $K'$ by field operations.
The established signs guarantee all needed denominators. No individual
conjugate of $\gamma$ is assumed available at this stage.

Construct $a=\sqrt A$, $c=\sqrt q$, $\rho=\sqrt{R_0}$ and $b=m/(4a)$.
The line $X=a$ and circle of center $(b,c)$ and radius $\rho$ are now
specified by constructed data. By Lemma~\ref{lem:signed} and $H_x(u)=0$,
$a_5=-mh-2h^2$. Since $D>0$, $P<0$, we have $\lambda>0$ and $h>0$.
The sextic in \eqref{eq:sextic} therefore meets every hypothesis of
Lemma~\ref{lem:reconstruct}.

Each real number $z_i=\lambda\beta_i$ is a root of this sextic.
Select the placements corresponding to $z_1$ and $z_3$, in successive
operations with the same prepared line and circle. Field operations
recover $\beta_1=z_1/\lambda$ and $\beta_3=z_3/\lambda$. Because
$P\ne0$, every $g_i$ and $\beta_i$ is nonzero. The identity
\begin{equation}\label{eq:recover}
\frac P{\beta_1\beta_3}
=\frac{g_0g_1g_2g_3g_4}{(g_1g_2)(g_3g_4)}
=g_0=1+u\gamma
\end{equation}
recovers $\gamma=(P/(\beta_1\beta_3)-1)/u$, since $u<-1$.

Finally $x_0<x_1$ gives $\sigma(\gamma)\ne\gamma$, so $\gamma\notin K$.
The prime degree $[L:K]=5$ implies $K(\gamma)=L$ by the tower law.
Every element of $L$ is a polynomial in $\gamma$ with coefficients in
$K$ and is obtained by finitely many field operations. Each step used
one of the coordinate rules, so Lemma~\ref{lem:sound} supplies a finite
geometric RMC construction.
\end{proof}

\begin{remark}
The recovery \eqref{eq:recover} avoids any need to prove that a single
adjacent product generates a compositum after the auxiliary and square
roots are adjoined. It also avoids requiring distinct adjacent-product
conjugates. The result is a finite existential construction, rather than
an optimized or numerically stable construction algorithm.
\end{remark}

\section{An exact conductor-25 example}\label{sec:example}

Set $t=(\sqrt5-1)/2$, $K=\Q(t)$, and
$\alpha=2\cos(2\pi/25)$. Then
\begin{equation}\label{eq:alpha}
t^2=1-t,\qquad \alpha^5-5\alpha^3+5\alpha-t=0.
\end{equation}
The real cyclotomic field $L=\Q(\alpha)$ has degree ten and contains
the real conductor-five field $K$. Thus $L/K$ is cyclic of degree five.
The angle-exponent action of 6 generates its Galois group, and reduction
of $2\cos(6\theta)$ by \eqref{eq:alpha} gives
\[
\sigma(\alpha)=-\alpha^4+4\alpha^2+t\alpha-2.
\]
Choose
\begin{equation}\label{eq:gamma25}
\gamma=\frac{-4-2\alpha+19\alpha^2-5\alpha^4}{4}.
\end{equation}
Exact arithmetic in the quotient defined by \eqref{eq:alpha} gives the
following equation for the auxiliary root, equal to $8192H_x(u)=0$:
\begin{equation}\label{eq:cubic25}
\begin{split}
0={}&(2991485+1857240t)u^3+(6475990+3876900t)u^2\\
&+(6055680+3653008t)u+(3136640+1917440t).
\end{split}
\end{equation}
There is a root in $I=(-119356/100000,-119355/100000)$, since exact
enclosures for the unscaled residual give
\begin{align*}
-21/5000&\le H_x(-119356/100000)\le-41/10000<0,\\
0<17/10000&\le H_x(-119355/100000)\le9/5000.
\end{align*}
Table~\ref{tab:example} records bounds valid throughout the closed
auxiliary interval. Also $P<0$ there, so the geometry is admissible.

\begin{table}[ht]
\centering
\caption{Exact conductor-25 geometric enclosures.}\label{tab:example}
\begin{tabular}{lrr}
\toprule Quantity&Lower bound&Upper bound\\\midrule
$D$&$101461/10000$&$50733/5000$\\
$\lambda$&$1971/10000$&$493/2500$\\
$h$&$1151/5000$&$2303/10000$\\
$q$&$77/625$&$309/2500$\\
$A$&$89/1250$&$763/10000$\\
$R_0$&$1389/2500$&$1143/2000$\\\bottomrule
\end{tabular}
\end{table}

For the real-embedding certificate, $\sqrt5$ is enclosed by consecutive
rationals with denominator $10^{30}$, beginning at
$2236067977499789696409173668731/10^{30}$. Squaring verifies this
enclosure. Similarly, $\alpha$ lies between
$1937166322257262238980336750929/10^{30}$ and that rational plus
$10^{-30}$. The quintic in \eqref{eq:alpha} has opposite signs at the
endpoints throughout the $t$ enclosure and has positive derivative
there. Its largest real root is $\alpha$: the roots are
$2\cos(2\pi r/25)$ for $r=1,6,11,16,21$, and the largest is the
$r=1$ value. The isolating interval is above $19/10$; the quintic is
strictly increasing on $[19/10,\infty)$ and changes sign before 2,
identifying that root.

The exact program verifies the automorphism, its order, and descent of
the coefficients of $P(U)$ and $e_j(U)$ to $K[U]$. Their exported values
specify the geometry without supplying $\alpha$ as a construction input.
After constructing $u$ and the parameters, obtain $\beta_1,\beta_3$
by the selected sextic placements and use \eqref{eq:recover} to recover
$\gamma$. A fixed exact identity then yields
\begin{equation}\label{eq:recover25}
\alpha=\frac{\sum_{j=0}^4(A_j+B_jt)\gamma^j}{45170808970351},
\end{equation}
with the integer coefficients in Table~\ref{tab:recover}. It is verified
by substituting \eqref{eq:gamma25} and reducing by \eqref{eq:alpha}.
Thus its use requires only field operations on already available values.
Constructing $\alpha/2$ and the corresponding positive sine gives a point
at central angle $2\pi/25$ on the unit circle and hence the regular
25-gon.

\begin{table}[ht]
\centering
\caption{Coefficients for the exact recovery \eqref{eq:recover25}.}\label{tab:recover}
\begin{tabular}{rrr}
\toprule $j$&$A_j$&$B_j$\\\midrule
0&$-82999543578141$&$-58480846421497$\\
1&$-13850695515680$&$-108800134011388$\\
2&$35174826943120$&$51032088536896$\\
3&$-11192002202560$&$42345158096896$\\
4&$3353265729280$&$-15629538309120$\\\bottomrule
\end{tabular}
\end{table}

\Needspace{11\baselineskip}
\section{A cyclotomic consequence}\label{sec:cyclotomic}

\begin{corollary}\label{cor:polygons}
Let $n\ge3$. If every prime factor of $\varphi(n)$ belongs to
$\{2,3,5\}$, then the regular $n$-gon is RMC-constructible.
\end{corollary}
\begin{proof}
The real cyclotomic extension
$\Q(\zeta_n+\zeta_n^{-1})/\Q$ is Galois and abelian of degree
$\varphi(n)/2$. Its Galois group admits a subgroup chain with successive
quotients cyclic of orders two, three, or five. Galois correspondence
gives a real field tower of those cyclic degrees. Quadratic steps are
compass constructions, cubic steps follow from Lemma~\ref{lem:cubic},
and quintic steps follow from Theorem~\ref{thm:main}. Thus
$\cos(2\pi/n)$ is constructible. A square root supplies
$\sin(2\pi/n)>0$, and repeated rotations construct the polygon.
\end{proof}

Examples include $n=25,31,41,61,101,125$. This is a sufficient condition;
no necessity assertion is needed here.

\section{Formal verification and reproducibility}\label{sec:verification}

The formalization uses Lean~4 \cite{Lean4} and Mathlib \cite{Mathlib}.

The finite coordinate predicate \code{Constructible} is defined in the
accompanying source
\begin{center}
\code{CyclicQuinticRMC_SET_FIX_V2.lean}.
\end{center}
Its main theorem is
\begin{center}
\code{RMCQuintic.real_galois_quintic_constructible}.
\end{center}
It assumes fields $K,L$, a finite-dimensional Galois $K$-algebra $L$ of
dimension five, and an embedding $L\to\R$, and concludes constructibility
of the embedded image of every element of $L$ from the embedded base
field. There is no general cubic-, quintic-, or sextic-root constructor.
The conditional closure theorem is discharged by the least set generated
by the explicit finite rules, rather than left as an extra hypothesis.

The pinned environment is Lean 4.34.0 with Mathlib v4.34.0, commit
\begin{center}\code{5ed2965256430c3649e86755f9576b54eca72435}.
\end{center}
The terminal dependency check reports
\begin{center}\code{[propext, Classical.choice, Quot.sound]}.
\end{center}
These are the standard foundational dependencies. In particular no
\code{sorryAx} or problem-specific axiom occurs in the dependency list.
The source was also checked with Lean 4.36.0-rc1. Table~\ref{tab:formal}
locates the main proof components. The geometric soundness
Lemma~\ref{lem:sound}, the conductor-25 identification and explicit
calculation, and Corollary~\ref{cor:polygons} are written arguments;
they are not claimed as additional formal theorems.

\begin{table}[ht]
\centering
\caption{Selected declarations in the formal proof.}\label{tab:formal}
\begin{tabular}{>{\raggedright\arraybackslash}p{.31\linewidth}>{\raggedright\arraybackslash}p{.60\linewidth}}
\toprule Component&Declarations\\\midrule
Conjugate density&\code{rationalConjugateTuple_hits_open} in namespace \code{CyclicQuintic}\\
Universal identity&\code{residual_identity}, \code{signed_obstruction}\\
Admissible specialization&\code{template_good}, \code{choose_good_conjugates}\\
Coefficient descent&\code{invariant_eval_mem}, \code{P_aux_mem}, the \code{e1_aux_mem}--\code{e5_aux_mem} lemmas\\
Geometric witnesses&\code{verge_root_mem}, \code{sextic_root_mem}, \code{cubic_mem}\\
Generator recovery&\code{adjacent_recover}, \code{generator_of_moved}\\
Final closure&\code{cyclic_quintic_closed}, \code{generated_closed}\\\bottomrule
\end{tabular}
\end{table}

The standard-library Python files \code{audit.py} and \code{example25.py}
provide exact rational-polynomial and interval checks. The first checks
the universal identities, cyclic invariance, geometric coefficients,
and template signs; the second performs exact quotient-field arithmetic
and certifies the conductor-25 intervals and descended input polynomials.
The file \code{geometry_audit_checks.py} provides 18 further symbolic
coefficient checks of the reconstruction and cubic identities.
The file \code{recovery25_check.py} verifies the integer recovery table
directly by exact quotient-field arithmetic. These checks do not use
sampled floating-point placements. Their scope is finite algebra and
inequalities; the geometric interpretation is the written proof above.

The supplementary package contains the complete Lean source, a pinned
Lake project, these Python files, their exact certificates, and build
instructions. In the pinned Lake environment run
\begin{verbatim}
lake env lean CyclicQuinticRMC_SET_FIX_V2.lean
\end{verbatim}
and compare the final dependency list. The source SHA-256 digest is
\begin{center}
\code{db36d978c9c88796885ebcbed464310d5ad2}\\
\code{c69c3ff6208e2902b5a7eb5bf9e7}.
\end{center}

\subsection*{Disclosure of generative AI assistance}
OpenAI ChatGPT was used extensively in developing the mathematical
argument, preparing the Lean proof and exact computational checks, and
drafting and editing the manuscript. The formal-verification claims are
limited to the declarations and dependencies specified above. The
geometric interpretation and supplementary written consequences remain
subject to mathematical scrutiny in the usual way.

\begin{thebibliography}{9}
\bibitem{BS14}
E.~Benjamin and C.~Snyder,
\emph{On the construction of the regular hendecagon by marked ruler and
compass}, Math. Proc. Cambridge Philos. Soc. \textbf{156} (2014),
409--424. \url{https://doi.org/10.1017/S0305004113000753}.

\bibitem{BS16}
E.~Benjamin and C.~Snyder,
\emph{Some fifth roots that are constructible by marked ruler and
compass}, Rocky Mountain J. Math. \textbf{46} (2016), 749--767.
\url{https://doi.org/10.1216/RMJ-2016-46-3-749}.

\bibitem{Lean4}
L.~de Moura and S.~Ullrich,
\emph{The Lean 4 theorem prover and programming language},
Automated Deduction--CADE 28, Lecture Notes in Comput. Sci.
\textbf{12699}, Springer, 2021, 625--635.
\url{https://doi.org/10.1007/978-3-030-79876-5_37}.

\bibitem{Mathlib}
The mathlib Community,
\emph{The Lean mathematical library}, Proceedings of the 9th ACM SIGPLAN
International Conference on Certified Programs and Proofs, ACM, 2020.
\url{https://doi.org/10.1145/3372885.3373824}.
\end{thebibliography}
\end{document}
